% !TEX program = platex
% 数学ノート集・共通テンプレート v2（積分論2026年度第1部の版面を参考）
% A5判。章内で定義・定理・例・注意等の番号を共通化。
% コンパイル： platex -interaction=nonstopmode math_note_template_v2.tex (2回)
%              dvipdfmx math_note_template_v2.dvi
\documentclass[a5paper,10pt,onecolumn]{jsarticle}
\usepackage[dvipdfmx,top=22mm,bottom=20mm,inner=19mm,outer=18mm,headheight=13pt,headsep=6mm]{geometry}
\usepackage{amsmath,amssymb,amsthm,mathtools}
\usepackage{enumitem,fancyhdr,titlesec}
\usepackage[dvipdfmx,bookmarks=true,colorlinks=false]{hyperref}
\usepackage{pxjahyper}

% ----- 差し替える文書情報 -----
\newcommand{\NoteTitle}{Lorentz幾何と一般相対論}
\newcommand{\NoteSeries}{数学ノート集}

% ----- 文字・組版 -----
% 参考PDFの質素な明朝＋Computer Modernの印象を引き継ぎながら，
% A5判で窮屈にならないよう本文の行間を少し確保する．
\linespread{1.13}
\setlength{\parindent}{1zw}
\setlength{\parskip}{0pt}
\setlength{\abovedisplayskip}{6pt plus 2pt minus 2pt}
\setlength{\belowdisplayskip}{6pt plus 2pt minus 2pt}
\setlength{\abovedisplayshortskip}{3pt plus 2pt minus 1pt}
\setlength{\belowdisplayshortskip}{4pt plus 2pt minus 1pt}
\setlist[enumerate]{topsep=3pt,itemsep=1pt,parsep=0pt,leftmargin=2.4zw}
\titleformat{\section}{\normalfont\normalsize\bfseries}{\thesection}{0.8zw}{}
\titlespacing*{\section}{0pt}{2.0ex plus .4ex}{.9ex}
\titleformat{\subsection}{\normalfont\normalsize\bfseries}{\thesubsection}{0.7zw}{}
\titlespacing*{\subsection}{0pt}{1.5ex plus .3ex}{.6ex}
\raggedbottom

% ----- 柱とページ番号：ページ上部、罫線なし -----
\fancypagestyle{mathnote}{%
  \fancyhf{}%
  \fancyhead[L]{\footnotesize \NoteSeries\ （\NoteTitle）}%
  \fancyhead[R]{\footnotesize \thepage}%
  \renewcommand{\headrulewidth}{0pt}%
  \renewcommand{\footrulewidth}{0pt}%
}
\pagestyle{mathnote}

% ----- 節内通し番号の定義・定理群 -----
\newtheoremstyle{mathnotetheorem}
 {4.5pt}{4.5pt}{\normalfont}{}%
 {\bfseries}{}{0.65zw}%
 {\thmname{#1}\thmnumber{ #2}\thmnote{（#3）}}
\theoremstyle{mathnotetheorem}
\newtheorem{theorem}{定理}[section]
\newtheorem{definition}[theorem]{定義}
\newtheorem{lemma}[theorem]{補題}
\newtheorem{proposition}[theorem]{命題}
\newtheorem{corollary}[theorem]{系}
\newtheorem{example}[theorem]{例}
\newtheorem{remark}[theorem]{注意}
\renewcommand{\proofname}{証明}
\renewcommand{\qedsymbol}{$\square$}
\makeatletter
\renewenvironment{proof}[1][\proofname]{%
  \par\pushQED{\qed}\normalfont\topsep4\p@\@plus2\p@\relax
  \trivlist\item[\hskip\labelsep\bfseries #1]\ignorespaces
}{\popQED\endtrivlist\@endpefalse}
\makeatother

% ----- 使うノートで適宜拡張する記号 -----
\newcommand{\R}{\mathbb{R}}
\newcommand{\norm}[1]{\lVert #1\rVert}
\newcommand{\abs}[1]{\lvert #1\rvert}

\begin{document}
\section*{到達目標と前提}
Lorentz内積から時間的・光的・空間的な変位を判別し，
世界線の固有時を計算できるようになることを主目標とする．
特に平坦な時空では固定した両端点を結ぶ慣性運動が
固有時を最大にすることを数学的に示す．
前提は多変数微分積分と初等的な線形代数である．

\section{Lorentz内積と因果構造}
$\mathbb R^{1+n}$ の座標を $(t,x_1,\ldots,x_n)$ とし
\[
 g(v,w)=-v_0w_0+\sum_{i=1}^nv_iw_i
\]
をLorentz内積とする．
\begin{definition}
非零ベクトル $v$ が時間的・光的・空間的であるとは，それぞれ
$g(v,v)<0$，$g(v,v)=0$，$g(v,v)>0$ を満たすことをいう．
未来向きであるとは時間成分 $v_0>0$ とする．
\end{definition}
\begin{proposition}
二つの未来向き時間的ベクトル $u,v$ の和は未来向き時間的である．
\end{proposition}
\begin{proof}
$u_0>\norm{\mathbf u}$，$v_0>\norm{\mathbf v}$ だから
$u_0+v_0>\norm{\mathbf u}+\norm{\mathbf v}\ge\norm{\mathbf u+\mathbf v}$．
よって $g(u+v,u+v)<0$ であり時間成分も正である．
\end{proof}
\section{Lorentz変換}
$|v|<1$ に対し $\gamma=(1-v^2)^{-1/2}$ と置く．
\begin{theorem}
二次元の線形変換
\[
 \begin{pmatrix}t'\\x'\end{pmatrix}
 =\begin{pmatrix}\gamma&-\gamma v\\-\gamma v&\gamma\end{pmatrix}
 \begin{pmatrix}t\\x\end{pmatrix}
\]
は $-dt^2+dx^2$ を保存する．
\end{theorem}
\begin{proof}
$t'=\gamma(t-vx)$，$x'=\gamma(x-vt)$ を代入すると
$-(t')^2+(x')^2=\gamma^2(1-v^2)(-t^2+x^2)=-t^2+x^2$．
極化公式により双線形形式も保存する．
\end{proof}
\section{曲がった時空の測地線}
滑らかなLorentz計量 $g_{ij}(x)$ を持つ多様体の局所座標で
$\Gamma^k_{ij}=\frac12g^{k\ell}(\partial_i g_{j\ell}
+\partial_jg_{i\ell}-\partial_\ell g_{ij})$ とする．
\begin{theorem}[測地線方程式]
$g(\dot\gamma,\dot\gamma)\ne0$ の定速曲線について，
エネルギー作用 $\frac12\int g_{ij}(x)\dot x^i\dot x^jdt$ の
固定端点変分の停留条件は
\[
 \ddot x^k+\Gamma^k_{ij}(x)\dot x^i\dot x^j=0
\]
である．
\end{theorem}
\begin{proof}
$L=\frac12g_{ij}\dot x^i\dot x^j$ のEuler--Lagrange方程式は
$\frac d{dt}(g_{kj}\dot x^j)-\frac12\partial_k g_{ij}\dot x^i\dot x^j=0$．
展開して逆行列 $g^{\ell k}$ を掛け，
$\dot x^i\dot x^j$ の対称性を用いると表示式になる．
\end{proof}
\begin{remark}
曲率テンソルとBianchi恒等式は後続の節で扱う．
Einstein方程式は数学的定義から証明される定理ではなく，
初期値問題の可解性は別の偏微分方程式論に属する．
\end{remark}


\section{Lorentz計量の接続と曲率}
\begin{theorem}[Levi--Civita接続]
滑らかなLorentz計量 $g$ を持つ多様体上には，
$\nabla_XY-\nabla_YX=[X,Y]$ と $\nabla g=0$ を満たす
接続 $\nabla$ がただ一つ存在する．局所座標での接続係数は
\[
 \Gamma^k_{ij}=\frac12g^{k\ell}
 (\partial_i g_{j\ell}+\partial_jg_{i\ell}-\partial_\ell g_{ij})
\]
で与えられる．
\end{theorem}
\begin{proof}
$\nabla g=0$ を用いて
$Xg(Y,Z)=g(\nabla_XY,Z)+g(Y,\nabla_XZ)$ のように展開し，
$X,Y,Z$ の循環置換で得られる式を加減する．
捩率零の条件で $\nabla_YX=\nabla_XY-[X,Y]$ を代入すると
\[
\begin{aligned}
2g(\nabla_XY,Z)={}&Xg(Y,Z)+Yg(X,Z)-Zg(X,Y)\\
 &-g(X,[Y,Z])+g(Y,[Z,X])+g(Z,[X,Y]).
\end{aligned}
\]
右辺は $g$ とLie括弧から一意に決まる．
計量が非退化だからこの式は $\nabla_XY$ を一意に定める．
右辺の $X,Y,Z$ に関する線形性とLeibniz則を直接確認すると，
得られた $\nabla$ は接続であり，式を引き算・足し算して
捩率零と計量両立性が成り立つことが確かめられる．
座標ベクトル場は互いに可換だから，$X=\partial_i,Y=\partial_j,Z=\partial_\ell$
を代入して逆計量を掛けると表示したChristoffel記号を得る．
\end{proof}
\begin{definition}
曲率作用素を
$R(X,Y)Z=\nabla_X\nabla_YZ-\nabla_Y\nabla_XZ-\nabla_{[X,Y]}Z$
と定める．座標成分では
\[
 R^\ell{}_{kij}=\partial_i\Gamma^\ell_{jk}
 -\partial_j\Gamma^\ell_{ik}
 +\Gamma^m_{jk}\Gamma^\ell_{im}
 -\Gamma^m_{ik}\Gamma^\ell_{jm}.
\]
Ricciテンソルを $\mathrm{Ric}_{kj}=R^i{}_{kij}$，
スカラー曲率を $S=g^{kj}\mathrm{Ric}_{kj}$，
Einsteinテンソルを $G=\mathrm{Ric}-\tfrac12Sg$ と定める．
\end{definition}
\begin{proposition}[Minkowski計量の平坦性]
$g=-dt^2+\sum_i dx_i^2$ では $R=0$ であり，
$\mathrm{Ric}=S=G=0$ である．
\end{proposition}
\begin{proof}
定数係数計量では全ての $\partial_i g_{jk}$ が零なので
全ての $\Gamma^k_{ij}=0$ である．曲率の座標公式へ代入すると
全ての成分が零であり，縮約したテンソルも零となる．
\end{proof}
\begin{remark}
Einstein方程式 $G+\Lambda g=8\pi T$ は，物質の応力エネルギーテンソル
$T$ と時空の計量を結ぶ物理法則であり，数学的な接続の定義のみから
「証明」されるものではない．その解の存在・一意性および
Bianchi恒等式の一般的な証明は，なお別の議論を要する．
\end{remark}

\section{Bianchi恒等式とEinsteinテンソルの発散}
\begin{lemma}[第一Bianchi恒等式]
Levi--Civita接続の曲率は任意のベクトル場 $X,Y,Z$ に対して
\[
 R(X,Y)Z+R(Y,Z)X+R(Z,X)Y=0
\]
を満たす．
\end{lemma}
\begin{proof}
左辺は各引数について $C^\infty(M)$-線形である．
点 $p$ の正規座標を取ると $\Gamma^\ell_{ij}(p)=0$ である．
曲率成分はその点で
$R^\ell{}_{kij}=\partial_i\Gamma^\ell_{jk}
-\partial_j\Gamma^\ell_{ik}$ となる．
さらに捩率零から $\Gamma^\ell_{ij}=\Gamma^\ell_{ji}$ である．
従って $R^\ell{}_{kij}+R^\ell{}_{ijk}+R^\ell{}_{jki}$ は
\[
 \partial_i\Gamma^\ell_{jk}-\partial_j\Gamma^\ell_{ik}
 +\partial_j\Gamma^\ell_{ki}-\partial_k\Gamma^\ell_{ji}
 +\partial_k\Gamma^\ell_{ij}-\partial_i\Gamma^\ell_{kj}=0.
\]
任意の $p$ でこの等式が成り立つから結論が従う．
\end{proof}
\begin{lemma}[第二Bianchi恒等式]
$\nabla$ を捩率零のLevi--Civita接続とし，曲率を
$R(X,Y)=[\nabla_X,\nabla_Y]-\nabla_{[X,Y]}$ とする．
このとき任意のベクトル場 $X,Y,Z$ に対し
\[
 (\nabla_XR)(Y,Z)+(\nabla_YR)(Z,X)
 +(\nabla_ZR)(X,Y)=0
\]
がEnd$(TM)$-値テンソルとして成立する．
\end{lemma}
\begin{proof}
両辺は $X,Y,Z$ について $C^\infty(M)$-線形なので，
点 $p$ で局所座標ベクトル場に代入して調べれば十分である．
接続の正規座標を $p$ の近傍で選べば $\Gamma^k_{ij}(p)=0$ である．
曲率成分は
\[
 R^\ell{}_{kij}=\partial_i\Gamma^\ell_{jk}
 -\partial_j\Gamma^\ell_{ik}
 +\Gamma^m_{jk}\Gamma^\ell_{im}
 -\Gamma^m_{ik}\Gamma^\ell_{jm}.
\]
$p$ では共変微分は通常の偏微分となる．
また右辺の二次項を偏微分すると，$p$ では
$\Gamma(p)=0$ により零になる．従って
$\nabla_aR^\ell{}_{kbc}+\nabla_bR^\ell{}_{kca}
 +\nabla_cR^\ell{}_{kab}$ は，
$\partial_a\partial_b\Gamma^\ell_{ck}$ などの
二階偏微分を交代和したものに等しい．
偏微分の交換可能性によってこの交代和は零となる．
$p$ は任意だから恒等式が大域的に成立する．
\end{proof}
\begin{theorem}[縮約Bianchi恒等式]
滑らかなLorentz計量 $g$ について
$G_{ab}=\operatorname{Ric}_{ab}-\frac12Sg_{ab}$ と置くと
\[
 \nabla^aG_{ab}=0
\]
が恒等的に成立する．
\end{theorem}
\begin{proof}
計量両立性 $\nabla g=0$ から曲率テンソルには
$R_{abcd}=-R_{bacd}=-R_{abdc}=R_{cdab}$ の対称性がある．
捩率零から生じる第一Bianchi恒等式も成立する．
ここでは $R_{abcd}=g(R(\partial_a,\partial_b)\partial_c,\partial_d)$，
$\operatorname{Ric}_{bd}=g^{ac}R_{abdc}$ と添字を定める．
第二Bianchi恒等式を
\[
 \nabla_a R_{bcde}+\nabla_bR_{cade}
 +\nabla_cR_{abde}=0
\]
と書き，$a,d$ を計量で縮約する．曲率の反対称性から
$g^{ad}R_{cade}=\operatorname{Ric}_{ce}$，
$g^{ad}R_{abde}=-\operatorname{Ric}_{be}$ なので
\[
 \nabla^dR_{bcde}=\nabla_c\operatorname{Ric}_{be}
 -\nabla_b\operatorname{Ric}_{ce}.
\]
さらに $b,e$ を縮約する．左辺は
$\nabla^d\operatorname{Ric}_{cd}$，右辺は
$\nabla_cS-\nabla^b\operatorname{Ric}_{cb}$ である．
従ってRicciテンソルの対称性を用いて
\[
 2\nabla^a\operatorname{Ric}_{ab}=\nabla_bS
\]
となる．
したがって $\nabla^ag_{ab}=0$ により
$\nabla^aG_{ab}=\nabla^a\operatorname{Ric}_{ab}
 -\frac12\nabla_bS=0$ である．
\end{proof}
\begin{remark}
$G+\Lambda g=8\pi T$ を採用すれば
$\nabla^aT_{ab}=0$ が数学的な整合条件として従う．
これは方程式を定義から導出したという意味ではない．
Lorentz計量についての幾何学的恒等式と，物理法則としての
Einstein方程式は論理的に区別される．
\end{remark}

\section{Einstein方程式の明示的な解：de Sitter時空}
四次元時空の座標を $(t,x^1,x^2,x^3)$ とし，
$a(t)>0$ に対してLorentz計量
\[
 g=-dt^2+a(t)^2\bigl((dx^1)^2+(dx^2)^2+(dx^3)^2\bigr)
\]
を考える．以下の計算では $i,j=1,2,3$ とする．
\begin{lemma}[平坦空間断面のRicciテンソル]
上の計量の非零Christoffel記号は
\[
 \Gamma^0_{ij}=a\dot a\,\delta_{ij},\qquad
 \Gamma^i_{0j}=\Gamma^i_{j0}=(\dot a/a)\delta^i_j
\]
であり，Ricciテンソルとスカラー曲率は
\[
 R_{00}=-3\ddot a/a,\quad R_{0i}=0,\quad
 R_{ij}=(a\ddot a+2\dot a^2)\delta_{ij},\qquad
 R=6\bigl(\ddot a/a+(\dot a/a)^2\bigr)
\]
を満たす．
\end{lemma}
\begin{proof}
$g_{00}=-1$，$g_{ij}=a^2\delta_{ij}$ と
$g^{00}=-1$，$g^{ij}=a^{-2}\delta^{ij}$ を
Levi--Civita接続の公式に代入すると表示したChristoffel記号を得る．
$H=\dot a/a$ と置く．Ricciテンソルの座標表示
\[
 R_{\mu\nu}=\partial_\alpha\Gamma^\alpha_{\mu\nu}
 -\partial_\nu\Gamma^\alpha_{\mu\alpha}
 +\Gamma^\alpha_{\mu\nu}\Gamma^\beta_{\alpha\beta}
 -\Gamma^\beta_{\mu\alpha}\Gamma^\alpha_{\nu\beta}
\]
により，$R_{00}=-3\dot H-3H^2=-3\ddot a/a$ を得る．
空間成分については
\[
 R_{ij}=(a\ddot a+\dot a^2)\delta_{ij}
 +3\dot a^2\delta_{ij}-2\dot a^2\delta_{ij}
 =(a\ddot a+2\dot a^2)\delta_{ij}
\]
となり，交差成分 $R_{0i}$ は零である．
$R=g^{\mu\nu}R_{\mu\nu}$ を計算すれば
スカラー曲率の表示式を得る．
\end{proof}
\begin{theorem}[正の宇宙項を持つ真空解]
$\Lambda>0$ とし $H_0=\sqrt{\Lambda/3}$ と置く．
$a(t)=e^{H_0t}$ と定めれば
\[
 \operatorname{Ric}(g)=\Lambda g,\qquad
 G_{\mu\nu}+\Lambda g_{\mu\nu}=0,
 \quad G_{\mu\nu}=R_{\mu\nu}-\frac12 Rg_{\mu\nu}
\]
が成立する．従って上の計量は宇宙項 $\Lambda$ を持つ真空Einstein方程式の
明示的なLorentz解である．
\end{theorem}
\begin{proof}
$\dot a=H_0a$，$\ddot a=H_0^2a$ だから前補題より
$R_{00}=-3H_0^2=-\Lambda=\Lambda g_{00}$，
$R_{ij}=3H_0^2a^2\delta_{ij}=\Lambda g_{ij}$ となり，
$R_{0i}=0=\Lambda g_{0i}$ である．よって
$\operatorname{Ric}(g)=\Lambda g$．
四次元なので $R=4\Lambda$ が従い，
$G=\Lambda g-\frac12(4\Lambda)g=-\Lambda g$ である．
\end{proof}
\begin{remark}
これは対称性の高い仮定の下で方程式を直接解いた例である．
任意のEinstein初期データから局所解を構成する一般Cauchy問題は
調和ゲージによる準線形双曲型方程式の理論を要し，本節では未証明である．
\end{remark}


\section{固有時を計算する：平坦時空の比較原理}
ここでは光速を $1$ とする単位系を用いる．
\begin{definition}
$T>0$，$x:[0,T]\to\mathbb R^n$ を
区分的に $C^1$ 級で，
$|\dot x(t)|<1$ を各滑らかな区間で満たす世界線とする．
その\emph{固有時}を
\[
 \tau(x)=\int_0^T\sqrt{1-|\dot x(t)|^2}\,dt
\]
と定める．積分は区分的に定義する．
\end{definition}
\begin{theorem}[直線世界線の固有時最大性]
$x(0)=0$，$x(T)=a$ とすると，$|a|<T$ であり
\[
 \tau(x)\le\sqrt{T^2-|a|^2}
\]
が成立する．等号が成り立つのは
$x(t)=ta/T$ のとき，かつそのときに限る．
\end{theorem}
\begin{proof}
単位開球上の $F(v)=\sqrt{1-|v|^2}$ を考える．
Hesse行列は
\[
 D^2F(v)=-\frac{I}{\sqrt{1-|v|^2}}
 -\frac{vv^T}{(1-|v|^2)^{3/2}}
\]
であり，任意の非零ベクトルに対して負定値となる．
従って $F$ は狭義凹関数である．
速度の平均は微積分学の基本定理より
\[
 \frac1T\int_0^T\dot x(t)\,dt=\frac aT.
\]
平均速度は開単位球に属し，Jensenの不等式から
\[
 \frac{\tau(x)}T
 \le F\left(\frac1T\int_0^T\dot x(t)\,dt\right)
 =\sqrt{1-\frac{|a|^2}{T^2}}.
\]
これが所望の評価である．狭義凹性から
等号なら速度はほとんど至る所一定であり，
区分的な連続性から各区間で $\dot x=a/T$，
端点条件から $x(t)=ta/T$ となる．
逆向きは直接代入で確かめられる．
\end{proof}
\begin{example}[往復する時計と静止する時計]
一次元で $a=0$ とし，$0<v<1$ とする．
$[0,T/2]$ では速度 $v$，
$[T/2,T]$ では速度 $-v$ の世界線は
出発点へ戻る．その固有時は
\[
 \tau=T\sqrt{1-v^2}<T.
\]
同じ出発点・到着点をもつ静止世界線 $x=0$ の
固有時は $T$ である．この不等式は
両端点が一致する場合の上の定理の具体例である．
\end{example}
\section*{到達確認問題と解答}
\noindent\textbf{問題．}
$T=5$，$a=3$ の一次元平坦時空で，
両端を結ぶ慣性世界線の速度と固有時を求めよ．

\medskip\noindent\textbf{解答．}
$x(t)=3t/5$ で速度 $3/5$．
固有時は
$\int_0^5\sqrt{1-9/25}\,dt=4$ である．
上の定理によって他の未来向き時間的世界線の固有時は
これを超えない．
\section*{射程}
上の最大性は平坦時空における固定端点問題である．
曲がった時空全体に同じ大域的結論を移すことはできず，
共役点や大域因果構造を別途考察する必要がある．

\end{document}
